% !TeX root = ../L_legacy_system_architecture.tex
% The complete trip chain: four layers, ordered by response time.
\begin{tikzpicture}[font=\scriptsize,x=1mm,y=1mm,every node/.style={outer sep=0pt}]
\tikzset{act/.style={text width=104mm,anchor=north west,inner sep=4pt,
                     rounded corners=1.6pt,font=\scriptsize,align=left},
         br/.style ={text width=96mm,anchor=north west,inner sep=3pt,
                     font=\tiny,align=left,draw=black!25,fill=white,
                     rounded corners=1.2pt}}

\node[pwr,text width=104mm,anchor=north west] (fault) at (0,0)
  {\textbf{Fault detected} --- arc, reflected power, vacuum, cooling, \dots};

% ---- Layer 1 -------------------------------------------------------------
\node[act,draw=cPwr!70,fill=cPwr!6,text=cPwr!85!black] (l1) at (10,-14)
  {\textbf{LAYER 1 --- Fast Interlock Chassis / AIM} \hfill hardware\\
   {\tiny B132 --- installed end-to-end timing not measured here}};
\node[br] (l1a) at (18,-28)
  {Arc / power detectors $\rightarrow$ \textbf{SCR ENABLE removed} (fiber $\rightarrow$ B514).
   Primary current is actually interrupted in \qtyrange{4}{8}{\milli\second} at the
  commutation event (published PEP-II result).};
\node[br] (l1b) at (18,-38)
  {Arc / power detectors $\rightarrow$ \textbf{CROWBAR fired} (fiber $\rightarrow$ B514).
  Published PEP-II conduction delay: $\approx\qty{10}{\micro\second}$.};
\node[br] (l1c) at (18,-47)
  {FISTAT fault word $\rightarrow$ AIM $\rightarrow$ VXI IOC.};

% ---- Layer 2 -------------------------------------------------------------
\node[act,draw=cPlc!70,fill=cPlc!7,text=cPlc!80!black] (l2) at (10,-58)
  {\textbf{LAYER 2 --- MPS interfaces} \hfill distinct paths\\
   {\tiny Monitors collector power, reflected power, vacuum, cooling.
    Per-fault actions and installed timing require verification:}};
\node[br] (l2a) at (18,-74)
  {RF MPS hardwired interface $\rightarrow$ Fast IC / AIM.
   Confirm SCR-inhibit and crowbar action for each fault class.};
\node[br] (l2b) at (18,-83)
  {RF MPS Remote I/O status $\rightarrow$ IOC alarms and sequencing;
   not a direct hardware RF cutoff.};
\node[br] (l2c) at (18,-92)
  {Independent SPEAR MPS shutoff $\rightarrow$ RF\_FAULT $\rightarrow$
   immediate RFP drive inhibit, without IOC polling.};

% ---- Layer 3 -------------------------------------------------------------
\node[act,draw=cPlc!70,fill=cPlc!7,text=cPlc!80!black] (l3) at (10,-103)
  {\textbf{LAYER 3 --- SLC-500 HVPS PLC} \hfill scan / timer dependent\\
  {\tiny Independent controller, not an output of the RF MPS relay.
    Monitors HVPS oil, crowbar, transformer arc, overvoltage.}};
\node[br] (l3a) at (18,-119)
  {Supervisory SCR enable relay de-energized $\rightarrow$ HVPS disabled.};
\node[br] (l3b) at (18,-128)
  {HVPS status $\rightarrow$ VXI IOC via Remote I/O.};

% ---- Layer 4 -------------------------------------------------------------
\node[act,draw=cCtrl!70,fill=cCtrl!6,text=cCtrl!85!black] (l4) at (10,-139)
  {\textbf{LAYER 4 --- EPICS / \texttt{rf\_states.st}} \hfill supervisory,
  event driven\\ {\tiny Sequencing and reporting, not fast hardware protection.}};
\node[br] (l4a) at (18,-153)
  {Fault file capture (\texttt{/dat/FAULT*\_N}).};
\node[br] (l4b) at (18,-162)
  {Station state $\rightarrow$ OFF (orderly HVPS shutdown via
  \texttt{s\_go\_off}; includes a nominal 5 s ramp wait when applicable).};
\node[br] (l4c) at (18,-171)
  {Operator alarm and notification.};

% ---- spine ---------------------------------------------------------------
% Stubs leave the rail at each layer box's own anchor height, so they are
% exactly horizontal regardless of how tall the box turns out.
\coordinate (rail) at (5,0);
\draw[bus] (5,-6) -- (5,-176);
\foreach \n in {l1,l2,l3,l4}{\draw[sig] (\n.west -| rail) -- (\n.west);}
\foreach \n in {l1a,l1b,l1c,l2a,l2b,l2c,l3a,l3b,l4a,l4b,l4c}
  {\draw[draw=black!35] (\n.west -| 14,0) ++(0,4) -- (\n.west -| 14,0) -- (\n.west);}
\end{tikzpicture}
